\documentclass[11pt,letterpaper]{report}

% ==============================================================================
% PAQUETES DE CONFIGURACIÓN BÁSICA Y TIPOGRAFÍA
% ==============================================================================
\usepackage[utf8]{inputenc}
\usepackage[T1]{fontenc}
\usepackage[english]{babel}
% Reconfiguracion de nombres en espanol (compatible universal)
\addto\captionsenglish{
  \renewcommand{\contentsname}{Tabla de Contenidos}
  \renewcommand{\chaptername}{Capítulo}
  \renewcommand{\tablename}{Tabla}
  \renewcommand{\figurename}{Figura}
}
\usepackage{lmodern}
\usepackage{microtype}

% Geometría de página equilibrada para manual técnico
\usepackage[left=2.5cm, right=2.5cm, top=2.8cm, bottom=2.8cm]{geometry}

% Paquetes matemáticos y químicos
\usepackage{amsmath, amssymb, amsfonts}
\usepackage{mathtools}

% Tablas profesionales y diseño
\usepackage{booktabs}
\usepackage{tabularx}
\usepackage{multirow}
\usepackage{array}
\usepackage{enumitem}

% Colores elegantes de ingeniería
\usepackage[table,xcdraw]{xcolor}
\definecolor{buapblue}{RGB}{0, 59, 112}      % Azul institucional clásico
\definecolor{deepcyan}{RGB}{0, 128, 128}     % Verde azulado petróleo
\definecolor{darkgray}{RGB}{45, 52, 54}      % Gris carbón de lectura
\definecolor{alertred}{RGB}{192, 57, 43}     % Rojo alerta metodológica
\definecolor{codegreen}{RGB}{39, 174, 96}    % Verde comentarios
\definecolor{codeblue}{RGB}{41, 128, 185}    % Azul palabras clave
\definecolor{codegray}{RGB}{127, 140, 141}   % Gris para números
\definecolor{codebg}{RGB}{248, 249, 250}     % Fondo claro para código
\definecolor{boxbg}{RGB}{240, 244, 248}      % Fondo cajas de texto
\definecolor{exercisebg}{RGB}{254, 249, 231} % Fondo suave para ejercicios

% Cajas estilizadas para notas, conceptos, alertas y ejercicios
\usepackage[most]{tcolorbox}

\newtcolorbox{conceptbox}[1]{
    colback=boxbg,
    colframe=buapblue,
    fonttitle=\bfseries\color{white},
    title=#1,
    sharp corners=downhill,
    arc=3mm,
    boxrule=1pt,
    left=4mm, right=4mm, top=3mm, bottom=3mm
}

\newtcolorbox{warningbox}[1]{
    colback=white,
    colframe=alertred,
    fonttitle=\bfseries\color{white},
    title=#1,
    sharp corners=downhill,
    arc=3mm,
    boxrule=1.2pt,
    left=4mm, right=4mm, top=3mm, bottom=3mm
}

\newtcolorbox{protocolbox}[1]{
    colback=white,
    colframe=deepcyan,
    fonttitle=\bfseries\color{white},
    title=#1,
    sharp corners=downhill,
    arc=3mm,
    boxrule=1pt,
    left=4mm, right=4mm, top=3mm, bottom=3mm
}

\newtcolorbox{exercisebox}[1]{
    colback=exercisebg,
    colframe=buapblue!80!black,
    fonttitle=\bfseries\color{white},
    title={Ejercicio Pedagógico: #1},
    sharp corners=downhill,
    arc=3mm,
    boxrule=1pt,
    left=4mm, right=4mm, top=3mm, bottom=3mm
}

% Formato de listados de código fuente con syntax highlighting
\usepackage{listings}
\lstdefinestyle{estilopython}{
    backgroundcolor=\color{codebg},
    commentstyle=\color{codegreen}\itshape,
    keywordstyle=\color{codeblue}\bfseries,
    numberstyle=\tiny\color{codegray},
    stringstyle=\color{alertred},
    basicstyle=\ttfamily\footnotesize\color{darkgray},
    breakatwhitespace=false,
    breaklines=true,
    captionpos=b,
    keepspaces=true,
    numbers=left,
    numbersep=8pt,
    showspaces=false,
    showstringspaces=false,
    showtabs=false,
    tabsize=4,
    frame=single,
    rulecolor=\color{codeblue!30}
}
\lstset{style=estilopython}

% Encabezados y pies de página estilizados
\usepackage{fancyhdr}
\pagestyle{fancy}
\fancyhf{}
\fancyhead[L]{\small\color{darkgray}\nouppercase{\leftmark}}
\fancyhead[R]{\small\bfseries\color{buapblue} Karol Paola Camacho López}
\fancyfoot[C]{\thepage}
\renewcommand{\headrulewidth}{0.6pt}
\renewcommand{\footrulewidth}{0.4pt}

% Hipervínculos interactivos
\usepackage{hyperref}
\hypersetup{
    colorlinks=true,
    linkcolor=buapblue,
    citecolor=deepcyan,
    urlcolor=codeblue,
    pdftitle={Manual Didáctico: Predicción de Entalpías de Combustión con Machine Learning - Karol Paola Camacho López},
    pdfauthor={Karol Paola Camacho López}
}

% Formato de títulos de capítulos y secciones
\usepackage{titlesec}
\titleformat{\chapter}[display]
  {\normalfont\bfseries\color{buapblue}}
  {\filleft\Huge\textbf{CAPÍTULO \thechapter}}{1ex}
  {\titlerule[1.5pt]\vspace{2ex}\LARGE}
\titlespacing*{\chapter}{0pt}{-20pt}{30pt}

% ==============================================================================
% DOCUMENTO PRINCIPAL
% ==============================================================================
\begin{document}

% ------------------------------------------------------------------------------
% PORTADA INSTITUCIONAL Y PERSONALIZADA
% ------------------------------------------------------------------------------
\begin{titlepage}
    \begin{center}
        \vspace*{0.3cm}
        {\LARGE \textbf{BENEMÉRITA UNIVERSIDAD AUTÓNOMA DE PUEBLA}}\\[0.3cm]
        {\Large \textbf{FACULTAD DE INGENIERÍA QUÍMICA}}\\[0.2cm]
        {\large Laboratorio de Termodinámica y Modelado Predictivo Molecular}\\[1.2cm]

        \begin{tcolorbox}[colback=buapblue!5,colframe=buapblue,arc=2mm,boxrule=1.5pt]
            \centering\vspace{0.3cm}
            {\huge \bfseries \color{buapblue} MANUAL DIDÁCTICO DE FORMACIÓN PROGRESIVA}\\[0.4cm]
            {\Large \bfseries Predicción de Entalpías Estándar de Combustión Mediante Machine Learning}\\[0.3cm]
            {\small \textit{De las Leyes Termodinámicas Clásicas al Modelado Predictivo de Alto Rendimiento}}\vspace{0.3cm}
        \end{tcolorbox}

        \vspace{1.2cm}

        {\large \textbf{Guía Integral de Aprendizaje, Ecuaciones, Código y Roadmap}}\\[0.3cm]
        {\normalsize Diseñado especialmente para la alumna:}\\[0.2cm]
        {\Large \bfseries \color{buapblue} KAROL PAOLA CAMACHO LÓPEZ}\\[0.2cm]
        {\normalsize Estudiante de Ingeniería Química $\cdot$ FIQ BUAP}\\[1.0cm]

        \vfill

        \begin{minipage}{0.88\textwidth}
            \centering
            \textbf{Enfoque Pedagógico:}\\
            Acompañamiento docente en Termodinámica Química, Ciencia de Datos y Machine Learning\\
            \vspace{0.3cm}
            \textbf{Objetivo Central de Aprendizaje:}\\
            Desarrollar y validar un modelo autónomo de regresión supervisada para la predicción de entalpías estándar de combustión ($\Delta H_c^\circ$) a partir de descriptores moleculares y fisicoquímicos, fundamentado en la literatura termodinámica rigurosa.\\
            \vspace{0.3cm}
            \textbf{Horizonte de Formación:}\\
            Fase Actual: Modelo Predictivo de Entalpías de Combustión (Septiembre -- Noviembre 2026)\\
            Meta Mayor: Transición al Modelado de Equilibrio Químico (10 de Diciembre de 2026)
        \end{minipage}

        \vspace{0.8cm}
        {\normalsize Heroica Puebla de Zaragoza, México $\cdot$ Septiembre 2026}
    \end{center}
\end{titlepage}

% ------------------------------------------------------------------------------
% TABLA DE CONTENIDOS
% ------------------------------------------------------------------------------
\tableofcontents
\newpage

% ------------------------------------------------------------------------------
% PREFACIO: CARTA DEL PROFESOR A KAROL PAOLA
% ------------------------------------------------------------------------------
\chapter*{Carta del Profesor a Karol Paola}
\addcontentsline{toc}{chapter}{Carta del Profesor a Karol Paola}

Estimada Karol Paola:

Bienvenida formalmente a este apasionante proyecto formativo en la Facultad de Ingeniería Química de nuestra universidad. Estás a punto de cruzar un puente fundamental para la ingeniería del siglo XXI: conectar los principios inmutables de la \textbf{Termodinámica Clásica} con la potencia predictiva del \textbf{Machine Learning}.

Sé muy bien que adentrarse en la inteligencia artificial y la ciencia de datos puede generar incertidumbre al principio. Hay una avalancha de términos nuevos: hiperparámetros, regresiones regularizadas, sobreajuste, matrices de dispersión, validación cruzada. Quiero que mantengas la calma y recuerdes siempre este principio pedagógico que guiará cada una de nuestras sesiones:
\begin{center}
    \textit{«El algoritmo es solo una calculadora avanzada de patrones estadísticos; la física y la química son las que dictan si el resultado tiene sentido o no.»}
\end{center}

Tu meta en esta primera etapa no es construir sistemas difusos ni memorizar decenas de algoritmos que no aplican a nuestro caso. Tu objetivo es claro, nítido y medible: \textbf{diseñar, entrenar, auditar y validar tu propio modelo predictivo de entalpía estándar de combustión ($\Delta H_c^\circ$)}, comprendiendo cada dato, cada ecuación termodinámica y cada línea de código.

No necesitas aprender todo Python de golpe ni dominar enciclopedias de estadística para comenzar. Trabajaremos bajo la metodología de \textbf{Aprendizaje Just-in-Time (JIT)}:
\begin{enumerate}[leftmargin=*]
    \item Estudiarás la ley física exacta del fenómeno (¿por qué la combustión libera calor?, ¿cómo mide la bomba calorimétrica?, ¿cómo se expresa matemáticamente?).
    \item Entenderás cómo traducir esa física a números que una computadora pueda procesar (masa molar, balance de oxígeno, conteo de heteroátomos).
    \item Aprenderás el comando específico de la librería adecuada (\texttt{pandas}, \texttt{numpy}, \texttt{chempy}, \texttt{scikit-learn}).
    \item Resolverás un ejercicio conceptual y lo aplicarás inmediatamente en tu propio cuaderno de trabajo.
\end{enumerate}

Este manual ha sido redactado como si estuviéramos sentados en el cubículo revisando tus dudas paso a paso. Úsalo como tu guía de cabecera. Pregunta todo lo que no quede claro, ejecuta cada celda de código por ti misma y disfruta el proceso de convertirte en una ingeniera química con dominio real de la ciencia de datos.

\vspace{0.6cm}
\begin{flushright}
    \textbf{Tu Profesor y Mentor de Proyecto}\\
    \textit{Laboratorio de Termodinámica y Modelado Molecular}\\
    Facultad de Ingeniería Química $\cdot$ BUAP
\end{flushright}

\newpage

% ------------------------------------------------------------------------------
% CAPÍTULO 1: TERMODINÁMICA FUNDAMENTAL Y LEYES DE LA COMBUSTIÓN
% ------------------------------------------------------------------------------
\chapter{Leyes Termodinámicas y Fenomenología de la Combustión}

Karol Paola, antes de solicitarle a un algoritmo de Machine Learning que prediga un número, debemos entender qué significa físicamente ese número. En este capítulo desglosaremos cada ley fundamental, cada variable de estado y cada ecuación que rige la liberación de energía en forma de calor.

\section{La Ley Cero: El Concepto Físico de Temperatura}

\subsection{Enunciado Físico y Sentido Experimental}
Si dos sistemas termodinámicos $A$ y $B$ se encuentran individualmente en equilibrio térmico con un tercer sistema $C$ (el termómetro calorimétrico), entonces los sistemas $A$ y $B$ están en equilibrio térmico entre sí:
\begin{equation}
    (T_A = T_C) \;\land\; (T_B = T_C) \implies T_A = T_B
\end{equation}

\noindent\textbf{¿Por qué le importa esto a tu modelo predictivo?}  
Porque en termoquímica todas las propiedades estándar se reportan a una temperatura de referencia idéntica: \textbf{$T^\circ = 298.15\text{ K}$ ($25.00^\circ\text{C}$)}. Si comparáramos una entalpía de combustión medida a $25^\circ\text{C}$ con otra medida a $100^\circ\text{C}$, el modelo aprendería un patrón falso, pues estaría mezclando la energía de reacción con el calor sensible ($C_p \Delta T$) de las sustancias.

\section{La Primera Ley: Conservación de Energía y el Origen de la Entalpía}

\subsection{Energía Interna ($U$), Calor ($Q$) y Trabajo ($W$)}
La Primera Ley postula que la energía del universo es constante: no se crea ni se destruye, solo se transforma. Para un sistema cerrado en reposo:
\begin{equation}
    dU = \delta Q - \delta W
\end{equation}
donde $U$ es la \textbf{energía interna} (suma de la energía cinética traslacional, rotacional, vibracional y la energía potencial electrostática de los enlaces covalentes e interacciones intermoleculares).

El trabajo cuasiestático de expansión o compresión de volumen frente a una presión externa $P$ se define como:
\begin{equation}
    \delta W = P\,dV \quad \Longrightarrow \quad dU = \delta Q - P\,dV
\end{equation}

\subsection{Deducción Matemática Rigurosa de la Entalpía ($H = U + PV$)}
En la mayoría de los reactores industriales, mecheros y llamas abiertas, las reacciones químicas ocurren a \textbf{presión constante} (la presión atmosférica circundante, $P = \text{constante} \implies dP = 0$).

Si integramos la Primera Ley a presión constante:
\begin{equation}
    \Delta U = Q_p - P\,\Delta V \quad \Longrightarrow \quad Q_p = \Delta U + P\,\Delta V
\end{equation}
Para evitar tener que calcular el cambio de volumen $\Delta V$ en cada reacción química, los termodinámicos definieron una nueva propiedad compuesta llamada \textbf{Entalpía ($H$)}:
\begin{equation}
    H \equiv U + PV
\end{equation}
Diferenciando esta definición:
\begin{equation}
    dH = dU + d(PV) = dU + P\,dV + V\,dP
\end{equation}
Sustituyendo la Primera Ley ($dU = \delta Q - P\,dV$):
\begin{equation}
    dH = (\delta Q - P\,dV) + P\,dV + V\,dP = \delta Q + V\,dP
\end{equation}
Bajo condición \textbf{isobárica} ($dP = 0$):
\begin{equation}
    dH = \delta Q_p \quad \Longrightarrow \quad \Delta H = Q_p
\end{equation}

\begin{conceptbox}{Lección Clave del Profesor}
La entalpía ($\Delta H$) no es un concepto místico: es simplemente el \textbf{calor absorbido o liberado por una reacción química cuando se lleva a cabo a presión constante}.
\begin{itemize}
    \item Si $\Delta H < 0$: El sistema libera calor hacia el entorno $\longrightarrow$ \textbf{Reacción Exotérmica} (Combustión).
    \item Si $\Delta H > 0$: El sistema absorbe calor del entorno $\longrightarrow$ \textbf{Reacción Endotérmica}.
\end{itemize}
\end{conceptbox}

\begin{exercisebox}{1.1: Calor Isobárico vs Cambio de Energía Interna}
\textbf{Problema:} Se quema $1\text{ mol}$ de metano gaseoso ($CH_4$) a $P = 1.0\text{ bar}$ y $T = 298.15\text{ K}$ según la reacción:
$$CH\_{4(g)} + 2\,O\_{2(g)} \longrightarrow CO\_{2(g)} + 2\,H\_2O\_{(l)}$$
Experimentalmente se mide una variación de entalpía $\Delta H^\circ = -890.8\text{ kJ/mol}$. Calcula el calor intercambiado a presión constante ($Q_p$) y la variación de energía interna ($\Delta U^\circ$). Considera comportamiento de gas ideal para los gases y volumen despreciable para el agua líquida ($R = 8.314\text{ J/(mol}\cdot\text{K)}$).

\vspace{0.2cm}\noindent\textbf{Solución Paso a Paso:}
\begin{enumerate}
    \item \textbf{Cálculo de $Q_p$:} Como la presión es constante, por definición directa:
    $$Q\_p = \Delta H^\circ = -890.8\text{ kJ/mol}$$
    \item \textbf{Cálculo del cambio de moles gaseosos ($\Delta n_g$):}
    $$\Delta n\_g = n\_{\text{gases, productos}} - n\_{\text{gases, reactivos}} = 1\text{ mol } (CO\_2) - [1\text{ mol } (CH\_4) + 2\text{ mol } (O\_2)] = 1 - 3 = -2\text{ mol}$$
    \item \textbf{Relación entre $\Delta H$ y $\Delta U$:}
    $$\Delta H = \Delta U + \Delta(PV) \approx \Delta U + \Delta n\_g R T$$
    $$\Delta U = \Delta H - \Delta n\_g R T$$
    \item \textbf{Sustitución numérica con cuidado de unidades:}
    $$\Delta n\_g R T = (-2\text{ mol}) \cdot \left(8.314 \times 10^{-3}\text{ kJ/(mol}\cdot\text{K)}\right) \cdot (298.15\text{ K}) = -4.958\text{ kJ}$$
    $$\Delta U = -890.8\text{ kJ} - (-4.958\text{ kJ}) = -885.84\text{ kJ/mol}$$
\end{enumerate}
\textit{Conclusión pedagógica:} La energía liberada como calor a presión constante ($-890.8\text{ kJ}$) es mayor en magnitud que la variación de energía interna ($-885.8\text{ kJ}$) porque la atmósfera realizó un trabajo de compresión ($P\Delta V < 0$) sobre el sistema al reducirse el número de moles gaseosos.
\end{exercisebox}

\section{La Segunda Ley: Espontaneidad y la Energía Libre de Gibbs}

\subsection{Entropía ($S$) e Irreversibilidad}
La combustión completa de un hidrocarburo es un proceso macroscópicamente irreversible: jamás observamos que el dióxido de carbono y el agua líquida se recombinen espontáneamente a temperatura ambiente para regenerar gasolina y oxígeno molecular. La Segunda Ley establece que en cualquier proceso espontáneo en un sistema aislado, la entropía total del universo aumenta:
\begin{equation}
    \Delta S_{\text{universo}} = \Delta S_{\text{sistema}} + \Delta S_{\text{alrededores}} > 0
\end{equation}

\subsection{La Ecuación de Gibbs: $\Delta G^\circ = \Delta H^\circ - T\Delta S^\circ$}
Para un sistema cerrado a temperatura y presión constantes, el calor cedido a los alrededores es $-\Delta H_{\text{sistema}}$, por lo que $\Delta S_{\text{alrededores}} = -\frac{\Delta H_{\text{sistema}}}{T}$. Sustituyendo:
\begin{equation}
    \Delta S_{\text{universo}} = \Delta S_{\text{sistema}} - \frac{\Delta H_{\text{sistema}}}{T} > 0
\end{equation}
Multiplicando por $-T$:
\begin{equation}
    -T \Delta S_{\text{universo}} = \Delta H_{\text{sistema}} - T \Delta S_{\text{sistema}} < 0
\end{equation}
Se define la \textbf{Energía Libre de Gibbs ($G$)} como:
\begin{equation}
    G \equiv H - TS \quad \Longrightarrow \quad \Delta G^\circ = \Delta H^\circ - T\Delta S^\circ
\end{equation}
Un proceso es termodinámicamente espontáneo a $T$ y $P$ constantes si y solo si $\Delta G < 0$. Como las reacciones de combustión liberan una cantidad inmensa de calor ($\Delta H_c^\circ \ll 0$), el término entálpico domina abrumadoramente, haciendo que $\Delta G_c^\circ$ sea intensamente negativo ($\approx -800$ a $-10000\text{ kJ/mol}$).

\subsection{Conexión con el Equilibrio Químico ($K_{eq}$)}
En la Fase final de tu formación (hacia el 10 de diciembre de 2026), relacionarás la energía libre con la constante de equilibrio químico mediante la ecuación de van 't Hoff:
\begin{equation}
    \Delta G^\circ = -RT \ln K_{eq} \quad \Longrightarrow \quad K_{eq} = \exp\left(-\frac{\Delta G^\circ}{RT}\right)
\end{equation}
Al ser $\Delta G_c^\circ$ enormemente negativo, $K_{eq} \ggg 10^{50}$, lo que demuestra termodinámicamente que la combustión procede al 100\% hasta agotar el reactivo limitante.

\section{Estequiometría de Combustión Completa y Ley de Hess}

\subsection{La Ecuación Estequiométrica General}
Para cualquier molécula orgánica neutra formada por carbono, hidrógeno, oxígeno y nitrógeno con fórmula molecular $C_a H_b O_c N_d$:
\begin{equation}
    C_a H_b O_c N_d + \nu_{O_2} O_{2(g)} \longrightarrow a\, CO_{2(g)} + \frac{b}{2}\, H_2O_{(l)} + \frac{d}{2}\, N_{2(g)}
\end{equation}
Realizando el balance de materia elemental para el oxígeno:
\begin{equation}
    c + 2\,\nu_{O_2} = 2a + \frac{b}{2} \quad \Longrightarrow \quad \nu_{O_2} = a + \frac{b}{4} - \frac{c}{2}
\end{equation}

\begin{warningbox}{¡Regla de Oro en Termodinámica de Combustión!}
En las convenciones del NIST y la IUPAC a $298.15\text{ K}$, el agua resultante de la combustión estándar se define en \textbf{estado líquido} ($H_2O_{(l)}$). Esto corresponde al \textbf{Poder Calorífico Superior (PCS)}. En calderas industriales se suele considerar agua vapor ($H_2O_{(g)}$, Poder Calorífico Inferior, PCI). Para el modelo de Machine Learning, entrenaremos siempre con el estándar científico internacional: \textbf{agua líquida}.
\end{warningbox}

\subsection{Cálculo de Entalpía por la Ley de Hess}
La entalpía es una \textbf{función de estado}: su cambio neto depende únicamente del estado inicial y del estado final, no de la trayectoria seguida. Por tanto:
\begin{equation}
    \Delta H_c^\circ = \sum \nu_{\text{prod}} \Delta H_{f,\text{prod}}^\circ - \sum \nu_{\text{react}} \Delta H_{f,\text{react}}^\circ
\end{equation}
Para la combustión de $C_a H_b O_c N_d$:
\begin{equation}
    \Delta H_c^\circ = a \Delta H_f^\circ(CO_{2(g)}) + \frac{b}{2} \Delta H_f^\circ(H_2O_{(l)}) + \frac{d}{2} \Delta H_f^\circ(N_{2(g)}) - \Delta H_f^\circ(\text{mol\'ecula}) - \nu_{O_2} \Delta H_f^\circ(O_{2(g)})
\end{equation}
Recordando que por convención termodinámica $\Delta H_f^\circ$ de los elementos puros en su forma más estable es cero ($\Delta H_f^\circ(O_{2(g)}) = 0$, $\Delta H_f^\circ(N_{2(g)}) = 0$):
\begin{equation}
    \Delta H_c^\circ = a (-393.51\text{ kJ/mol}) + \frac{b}{2} (-285.83\text{ kJ/mol}) - \Delta H_f^\circ(\text{mol\'ecula})
\end{equation}

\begin{exercisebox}{1.2: Balance y Cálculo de Entalpía por Ley de Hess}
\textbf{Problema:} La entalpía estándar de formación del etanol líquido ($C_2H_5OH_{(l)}$) a $298.15\text{ K}$ es $\Delta H_f^\circ = -277.69\text{ kJ/mol}$.
\begin{enumerate}
    \item Escribe y balancea la reacción de combustión completa.
    \item Calcula teóricamente su entalpía estándar de combustión ($\Delta H_c^\circ$).
    \item Calcula la entalpía específica liberada por gramo de etanol ($M = 46.07\text{ g/mol}$).
\end{enumerate}

\vspace{0.2cm}\noindent\textbf{Solución Paso a Paso:}
\begin{enumerate}
    \item \textbf{Fórmula molecular del etanol:} $C_2 H_6 O_1 \implies a=2, b=6, c=1, d=0$.
    $$\nu\_{O\_2} = a + \frac{b}{4} - \frac{c}{2} = 2 + \frac{6}{4} - \frac{1}{2} = 2 + 1.5 - 0.5 = 3$$
    Reacción balanceada:
    $$C\_2H\_5OH\_{(l)} + 3\,O\_{2(g)} \longrightarrow 2\,CO\_{2(g)} + 3\,H\_2O\_{(l)}$$
    \item \textbf{Aplicación de la Ley de Hess:}
    $$\Delta H\_c^\circ = 2\cdot(-393.51) + 3\cdot(-285.83) - (-277.69)$$
    $$\Delta H\_c^\circ = -787.02 - 857.49 + 277.69 = -1366.82\text{ kJ/mol}$$
    \item \textbf{Entalpía específica (Poder calorífico por gramo):}
    $$\Delta h\_c^\circ = \frac{\Delta H\_c^\circ}{M} = \frac{-1366.82\text{ kJ/mol}}{46.07\text{ g/mol}} = -29.67\text{ kJ/g}$$
\end{enumerate}
\end{exercisebox}

\section{Calorimetría de Bomba de Mahler: El Origen Experimental de tus Datos}

La enorme mayoría de los datos tabulares que procesarás provienen históricamente de un \textbf{calorímetro de bomba de Mahler}:

\subsection{El Principio de Medición a Volumen Constante}
La bomba calorimétrica es un recipiente de acero inoxidable de paredes gruesas sumergido en una masa conocida de agua. Como la bomba está sellada herméticamente:
\begin{equation}
    dV = 0 \quad \Longrightarrow \quad W = \int P\,dV = 0
\end{equation}
Por la Primera Ley ($dU = \delta Q - \delta W$):
\begin{equation}
    \Delta U_c = Q_v
\end{equation}
El calor desprendido eleva la temperatura del agua y de los componentes del calorímetro:
\begin{equation}
    Q_v = - C_{\text{calor\'imetro}} \cdot \Delta T_{\text{corregido}}
\end{equation}
donde $C_{\text{calor\'imetro}}$ (la capacidad calorífica efectiva o equivalente en agua del equipo) se calibra previamente quemando pastillas de ácido benzoico patrón de alta pureza ($\Delta_c u = -26434\text{ J/g}$).

\subsection{Conversión de $\Delta U_c$ a $\Delta H_c^\circ$}
Como la bomba mide $\Delta U_c$ y no $\Delta H_c$, se aplica la corrección termodinámica:
\begin{equation}
    \Delta H_c^\circ = \Delta U_c + \Delta n_g R T + \Delta H_{\text{Washburn}}
\end{equation}
donde $\Delta H_{\text{Washburn}}$ agrupa las correcciones por:
\begin{itemize}[leftmargin=*]
    \item La no idealidad de los gases reales a $30\text{ bar}$ de presión de oxígeno inicial.
    \item La energía liberada por el alambre de ignición de platino o níquel-cromo.
    \item El calor de formación y disolución de trazas de ácido nítrico ($HNO_3$) formado por nitrógeno atmosférico residual.
\end{itemize}

\section{¿Por qué Requerimos Machine Learning en Lugar de Medir Todo?}
Podrías preguntarte: si la bomba calorimétrica es tan exacta, ¿por qué no medimos experimentalmente todas las moléculas y nos olvidamos del Machine Learning?
\begin{enumerate}[leftmargin=*]
    \item \textbf{Demanda de Masa:} Un calorímetro de bomba requiere entre $0.5\text{ g}$ y $1.0\text{ g}$ de sustancia pura. La síntesis orgánica preliminar de fármacos o catalizadores novedosos rinde con frecuencia menos de $10\text{ mg}$.
    \item \textbf{Costo y Tiempo:} Una sola corrida calorimétrica rigurosa con calibración, repeticiones y correcciones de Washburn toma de 2 a 4 días de trabajo intensivo de laboratorio.
    \item \textbf{Seguridad Operativa:} Compuestos altamente reactivos, peróxidos, nitratos inestables o sustancias químicas tóxicas no pueden someterse a combustión a $30\text{ bar}$ sin riesgo de daño estructural o peligro para el experimentador.
\end{enumerate}

\newpage

% ------------------------------------------------------------------------------
% CAPÍTULO 2: CONFIGURACIÓN DEL LABORATORIO DIGITAL
% ------------------------------------------------------------------------------
\chapter{Configuración del Laboratorio Digital: PyCharm, Entornos y Git}

Karol Paola, antes de escribir una sola línea de código químico, debemos asegurar que tu espacio digital sea reproducible, ordenado y profesional. Un modelo científico no sirve si solo funciona en una computadora específica.

\section{Estructura Estándar del Repositorio}
Abre tu explorador de archivos y confirma que la carpeta raíz del proyecto contenga exactamente la siguiente arquitectura de directorios:

\begin{tcolorbox}[colback=codebg,colframe=darkgray,title=Estructura Estandarizada del Proyecto]
\begin{verbatim}
Modelo Entalpias de Combustion/
|-- data/
|   |-- raw/                  <-- Datos originales intactos (solo lectura)
|   |   `-- base_entalpias.csv
|   `-- processed/            <-- Tablas limpias con variables calculadas
|-- notebooks/
|   |-- 01_exploracion_inicial.ipynb  <-- Tu cuaderno del primer viernes
|   `-- 02_limpieza_y_eda.ipynb
|-- src/                      <-- Modulos Python reutilizables (.py)
|   |-- __init__.py
|   `-- funciones_quimicas.py <-- Balance de materia y utilidades
|-- config/                   <-- Configuraciones del proyecto
|   |-- matriz_actividades.json
|   `-- MATRIZ_ACTIVIDADES.md
|-- reportes/                 <-- Documentos tecnicos y este manual en LaTeX
|   `-- MANUAL_MAESTRO_PREPARACION.tex
|-- .gitignore                <-- Archivos excluidos de Git (.venv, cache)
|-- README.md                 <-- Descripcion oficial del repositorio
|-- KANBAN.md                 <-- Tablero de actividades sincronizado
|-- ROADMAP.md                <-- Diagrama cronologico de entrega
`-- requirements.txt          <-- Lista formal de dependencias y versiones
\end{verbatim}
\end{tcolorbox}

\section{Paso a Paso: Configurar el Entorno Virtual en PyCharm}

\begin{protocolbox}{Protocolo 2.1: Entorno Virtual Aislado (.venv)}
\begin{enumerate}[leftmargin=*]
    \item \textbf{Abrir la Carpeta:} Inicia PyCharm, selecciona \textbf{Open} y abre la carpeta \texttt{Modelo Entalpias de Combustion}.
    \item \textbf{Configurar Intérprete:} Ve a \textbf{File $\longrightarrow$ Settings} (en Linux/Windows) o \textbf{PyCharm $\longrightarrow$ Preferences} (en macOS).
    \item Despliega \textbf{Project: Modelo Entalpias...} y haz clic en \textbf{Python Interpreter}.
    \item Haz clic en \textbf{Add Interpreter $\longrightarrow$ Add Local Interpreter...}.
    \item Selecciona \textbf{Virtualenv Environment} $\longrightarrow$ \textbf{New Environment}.
    \item Verifica que la carpeta de destino termine en \texttt{.venv} y que el intérprete base sea Python 3.10 o superior. Haz clic en \textbf{OK}.
\end{enumerate}
\end{protocolbox}

\section{Instalación de Dependencias desde la Terminal Integrada}
Abre la pestaña \textbf{Terminal} en la parte inferior de PyCharm. Asegúrate de que el cursor comience con \texttt{(.venv)}. Ejecuta exactamente:

\begin{lstlisting}[language=bash]
# 1. Actualizar el instalador de paquetes pip
pip install --upgrade pip

# 2. Instalar el nucleo interactivo para cuadernos Jupyter
pip install jupyter jupyterlab ipykernel

# 3. Instalar la suite cientifica de calculo y Machine Learning
pip install numpy pandas chempy matplotlib seaborn scikit-learn
\end{lstlisting}

\section{Manejo Estricto de Rutas Relativas (Regla de Oro)}
Jamás escribas rutas absolutas con letras de disco o carpetas personales como:  
\texttt{pd.read\_csv("C:/Users/Karol/Documents/.../datos.csv")}. Si lo haces, el código fallará inmediatamente cuando se ejecute en la computadora del profesor.

Utiliza siempre \texttt{os.path.join()} o \texttt{pathlib.Path} basándote en la raíz del proyecto:
\begin{lstlisting}[language=Python]
import os
import pandas as pd

# Construir la ruta relativa desde la carpeta notebooks/ hacia data/raw/
ruta_datos = os.path.join("..", "data", "raw", "base_entalpias.csv")
df = pd.read_csv(ruta_datos)
print(f"Dimensiones del dataset cargado: {df.shape[0]} filas, {df.shape[1]} columnas.")
\end{lstlisting}

\newpage

% ------------------------------------------------------------------------------
% CAPÍTULO 3: NORMA DE ESTANDARIZACIÓN Y TRAZABILIDAD DOCUMENTAL
% ------------------------------------------------------------------------------
\chapter{Norma de Estandarización y Trazabilidad Documental}

Karol Paola, para que tu flujo de trabajo sea riguroso y para que las automatizaciones del proyecto reconozcan tus avances de forma determinista, cada entregable producido debe satisfacer la siguiente norma universal de nomenclatura.

\section{Estructura Universal de Nombres de Archivo}

\begin{tcolorbox}[colback=codebg,colframe=buapblue,arc=2mm,boxrule=1pt]
\centering\large\ttfamily
[CÓDIGO\_FASE]\_[ID\_ACTIVIDAD]\_[TIPO]\_[DESCRIPTOR]\_[VERSION]\_[ESTADO].[ext]
\end{tcolorbox}

\begin{table}[h!]
\centering
\small
\begin{tabularx}{\textwidth}{c l p{3.5cm} X}
\toprule
\textbf{Bloque} & \textbf{Campo} & \textbf{Formato / Valores} & \textbf{Función Técnica en el Sistema} \\
\midrule
\rowcolor{codebg}
\textbf{1} & \texttt{[CÓDIGO\_FASE]} & Prefijo de etapa (ej. \texttt{F0-ENV}, \texttt{F1-DAT}) & Ubica el entregable en la fase cronológica del Roadmap. \\
\addlinespace
\textbf{2} & \texttt{[ID\_ACTIVIDAD]} & 3 letras + 2 dígitos (ej. \texttt{ACT01}, \texttt{SIM01}) & Identificador unívoco vinculado a la tarjeta Kanban. \\
\rowcolor{codebg}
\textbf{3} & \texttt{[TIPO]} & 3 letras mayúsculas (\texttt{DAT}, \texttt{SCR}, \texttt{NBK}, \texttt{REP}) & Clasifica la naturaleza técnica del archivo. \\
\addlinespace
\textbf{4} & \texttt{[DESCRIPTOR]} & Texto en minúsculas en \texttt{snake\_case} & Síntesis temática del entregable (máximo 4 palabras). \\
\rowcolor{codebg}
\textbf{5} & \texttt{[VERSION]} & \texttt{vX.Y} (\texttt{v0.1}, \texttt{v1.0}, \texttt{v2.0}) & Versionado semántico ($X$: mayor; $Y$: menor). \\
\addlinespace
\textbf{6} & \texttt{[ESTADO]} & \texttt{DRAFT}, \texttt{REV}, \texttt{FINAL} & Determina si el archivo es borrador, revisión o cierre. \\
\rowcolor{codebg}
\textbf{Ext} & \texttt{.[ext]} & Extensión (\texttt{.py}, \texttt{.ipynb}, \texttt{.csv}, \texttt{.pdf}) & Formato nativo del software utilizado. \\
\bottomrule
\end{tabularx}
\caption{Parámetros normativos de nomenclatura estricta.}
\end{table}

\subsection{Prefijos Oficiales de Fase (\texttt{CÓDIGO\_FASE})}
\begin{itemize}[leftmargin=*]
    \item \textbf{\texttt{F0-ENV} (Entorno y Setup):} Infraestructura local, Git, entorno virtual \texttt{.venv} y ChemPy.
    \item \textbf{\texttt{F1-DAT} (Datos y Curaduría):} Adquisición de la base de datos de entalpías y depuración de duplicados.
    \item \textbf{\texttt{F2-EDA} (Exploración y Descriptores):} Análisis exploratorio de datos (EDA) y cálculo de variables moleculares.
    \item \textbf{\texttt{F3-MLB} (Machine Learning Base):} Ajuste de modelos lineales (OLS, Ridge) y Random Forest.
    \item \textbf{\texttt{F4-VAL} (Validación y Residuales):} Validación cruzada $k$-fold y diagnóstico de gráficos residuales.
    \item \textbf{\texttt{F5-REP} (Benchmarking de Literatura):} Comparación cuantitativa del modelo contra estándares reportados.
    \item \textbf{\texttt{F6-SIM} (Simulación de Procesos):} Balance termodinámico de combustión completa en DWSIM (\texttt{.dwsim}).
    \item \textbf{\texttt{F7-EQK} (Transición al Equilibrio):} Cálculo de energía de Gibbs y constantes de equilibrio $K_{eq}(T)$ para diciembre.
\end{itemize}

\subsection{Expresión Regular (Regex) de Validación}
Para que el motor de automatización valide formalmente un entregable antes de cerrar una tarea, el nombre de archivo debe coincidir con la siguiente regla lógica:
\begin{lstlisting}[language=bash]
^(F[0-7]-[A-Z]{3})_([A-Z]{3}\d{2})_([A-Z]{3})_([a-z0-9_]+)_(v\d+\.\d+)_(DRAFT|REV|FINAL)\.([a-zA-Z0-9]+)$
\end{lstlisting}

\newpage

% ------------------------------------------------------------------------------
% CAPÍTULO 4: LIBRERÍAS CIENTÍFICAS EXPLICADAS LÍNEA POR LÍNEA
% ------------------------------------------------------------------------------
\chapter{Librerías Científicas para Ingeniería Química: Guía Línea por Línea}

En este capítulo analizaremos las herramientas computacionales que utilizarás en tu día a día: NumPy, Pandas, ChemPy, Scikit-Learn y Matplotlib/Seaborn.

\section{NumPy: Arreglos Vectorizados de Alto Rendimiento}
NumPy implementa arreglos homogéneos contiguos en memoria C (\texttt{ndarray}). Permite procesar millones de cálculos químicos sin recurrir a bucles \texttt{for} en Python.

\begin{lstlisting}[language=Python]
import numpy as np

# 1. Definir arreglos vectorizados de entalpias molares (kJ/mol) y masas (g/mol)
entalpias_molares = np.array([-890.8, -1560.7, -2220.0, -2877.6]) # Metano, Etano, Propano, Butano
masas_molares = np.array([16.04, 30.07, 44.10, 58.12])

# 2. Operacion vectorizada directa: calculo de poder calorifico especifico (kJ/g)
entalpias_especificas = entalpias_molares / masas_molares
print("Entalpias especificas (kJ/g):", np.round(entalpias_especificas, 2))

# 3. Estadistica basica sobre el target
media_dH = np.mean(entalpias_molares)
desv_std = np.std(entalpias_molares)
print(f"Promedio: {media_dH:.2f} kJ/mol | Desv. Estandar: {desv_std:.2f} kJ/mol")

# 4. Mascara booleana para filtrar valores extremos
mascara_alta_energia = entalpias_molares < -2000.0
print("Compuestos con dH < -2000 kJ/mol:", entalpias_molares[mascara_alta_energia])
\end{lstlisting}

\section{Pandas: La Tabla de Datos del Ingeniero Químico}
Un \texttt{DataFrame} es una tabla interactiva en memoria donde las columnas son variables fisicoquímicas (\textit{features}) y las filas son moléculas individuales.

\begin{lstlisting}[language=Python]
import pandas as pd

# 1. Cargar la base de datos cruda
df = pd.read_csv('../data/raw/base_entalpias.csv')

# 2. Inspeccion inicial de dimensiones y tipos de columnas
print("Forma del dataset:", df.shape) # (num_filas, num_columnas)
print("\nTipos de datos y conteo de nulos:")
df.info()

# 3. Resumen estadistico de variables numericas (media, std, min, percentiles, max)
display(df.describe())

# 4. Auditoria y remocion de valores nulos (NaN)
# REGLA DOCENTE: El target JAMAS se inventa ni se rellena con promedios. Si falta dH, se descarta la fila.
num_nulos_target = df['dH_combustion'].isnull().sum()
print(f"Filas con entalpia nula: {num_nulos_target}")
df_limpio = df.dropna(subset=['dH_combustion']).copy()

# 5. Deteccion y remocion de duplicados moleculares basados en SMILES
df_limpio = df_limpio.drop_duplicates(subset=['SMILES'])

# 6. Ingenieria de Variables (Feature Engineering): creacion de nuevas columnas utiles
df_limpio['ratio_C_H'] = df_limpio['num_C'] / df_limpio['num_H']
\end{lstlisting}

\section{ChemPy: Estequiometría y Balance Automatizado}
ChemPy es una librería especializada de química computacional que permite extraer masas molares y balancear sistemas estequiométricos mediante álgebra lineal elemental.

\begin{lstlisting}[language=Python]
from chempy import Substance
from chempy.chemistry import balance_stoichiometry

# 1. Instanciar una molecula a partir de su formula quimica
tolueno = Substance.from_formula('C7H8')
print(f"Masa molar del tolueno: {tolueno.mass:.3f} g/mol")
print("Composicion elemental (numero atomico: conteo):", tolueno.composition)
# Salida: {6: 7, 1: 8} -> 7 atomos de Carbono (Z=6) y 8 atomos de Hidrogeno (Z=1)

# 2. Balancear automaticamente la reaccion de combustion completa con O2
reactivos = {'C7H8', 'O2'}
productos = {'CO2', 'H2O'}

reac_bal, prod_bal = balance_stoichiometry(reactivos, productos)
print("\nReaccion de combustion del tolueno balanceada:")
print(f"{dict(reac_bal)} --> {dict(prod_bal)}")
# Salida esperada: {'C7H8': 1, 'O2': 9} --> {'CO2': 7, 'H2O': 4}
\end{lstlisting}

\newpage

% ------------------------------------------------------------------------------
% CAPÍTULO 5: REPRESENTACIÓN MOLECULAR PARA EL MODELO PREDICTIVO
% ------------------------------------------------------------------------------
\chapter{Representación Molecular: De la Química a los Números}

Karol Paola, un algoritmo de Machine Learning no entiende de geometrías moleculares tridimensionales ni de resonancia electrónica; únicamente procesa matrices numéricas $\mathbf{X} \in \mathbb{R}^{n \times p}$. En este capítulo aprenderás cómo convertir una molécula orgánica en un vector de números con significado termodinámico.

\section{La Notación SMILES: Moléculas Escritas en Texto}
El estándar mundial para almacenar moléculas en bases de datos es el formato \textbf{SMILES} (\textit{Simplified Molecular Input Line Entry System}):
\begin{itemize}[leftmargin=*]
    \item \textbf{Átomos alifáticos:} Letras mayúsculas según la tabla periódica (\texttt{C, N, O, S, F, Cl, Br, I}). Los hidrógenos se asumen implícitos para satisfacer la valencia del átomo.
    \item \textbf{Aromaticidad:} Se indica con letras minúsculas. Por ejemplo, el benceno es \texttt{c1ccccc1}.
    \item \textbf{Tipos de enlace:} Enlace simple implícito; enlace doble con \texttt{=} (ej. etileno \texttt{C=C}); enlace triple con \texttt{\#} (ej. acetileno \texttt{C\#C}).
    \item \textbf{Ramificaciones:} Se colocan entre paréntesis justo después del átomo al que van enlazadas. Por ejemplo, el isobutano (2-metilpropano) se escribe \texttt{CC(C)C}.
\end{itemize}

\section{Vector de Descriptores Moleculares y Fisicoquímicos ($\mathbf{x}$)}

Para alimentar nuestro modelo de regresión, cada compuesto $i$ se representará como un vector fila:
\begin{equation}
    \mathbf{x}_i = \begin{bmatrix} M_i & n_{C,i} & n_{H,i} & n_{O,i} & n_{N,i} & \nu_{O_2,i} & \text{Ratio}_{C/H,i} & \text{EnlacesRotables}_i & \text{TPSA}_i \end{bmatrix}
\end{equation}

\noindent\textbf{Justificación termodinámica de cada descriptor:}
\begin{enumerate}[leftmargin=*]
    \item \textbf{Masa Molar ($M$):} Correlaciona directamente con el tamaño molecular y la cantidad total de electrones y núcleos presentes.
    \item \textbf{Conteo de Carbonos ($n_C$):} Cada carbono se oxida a $CO_2$, liberando aproximadamente $-393.5\text{ kJ/mol}$. Es el principal predictor positivo de la magnitud de la entalpía.
    \item \textbf{Conteo de Hidrógenos ($n_H$):} Cada dos hidrógenos forman un mol de $H_2O_{(l)}$, liberando $-285.8\text{ kJ/mol}$.
    \item \textbf{Conteo de Oxígenos ($n_O$):} Los oxígenos presentes en la molécula ya están enlazados parcialmente al carbono o hidrógeno; por tanto, reducen la cantidad de oxígeno externo que se debe consumir y disminuyen la energía neta liberada en la combustión.
    \item \textbf{Demanda Estequiométrica de Oxígeno ($\nu_{O_2} = a + b/4 - c/2$):} Sintetiza en una sola variable físico-química el potencial redox global de la molécula.
\end{enumerate}

\begin{exercisebox}{5.1: Construcción de un Vector de Descriptores}
\textbf{Problema:} Para las siguientes tres moléculas:
\begin{itemize}
    \item Molécula A: Metanol (\texttt{CO}, fórmula $CH_4O$)
    \item Molécula B: Acetona (\texttt{CC(=O)C}, fórmula $C_3H_6O$)
    \item Molécula C: Ácido acético (\texttt{CC(=O)O}, fórmula $C_2H_4O_2$)
\end{itemize}
Calcula manualmente su masa molar ($M$), su conteo atómico ($n_C, n_H, n_O$), su demanda estequiométrica de oxígeno ($\nu_{O_2}$) y su razón $C/H$.

\vspace{0.2cm}\noindent\textbf{Solución:}
\begin{center}
\small
\begin{tabular}{l c c c c c c}
\toprule
\textbf{Compuesto} & \textbf{Masa Molar ($M$)} & $n_C$ & $n_H$ & $n_O$ & $\nu_{O_2} = n_C + \frac{n_H}{4} - \frac{n_O}{2}$ & \textbf{Ratio $C/H$} \\
\midrule
Metanol ($CH_4O$) & $32.04\text{ g/mol}$ & 1 & 4 & 1 & $1 + 1 - 0.5 = \mathbf{1.5}$ & $0.25$ \\
Acetona ($C_3H_6O$) & $58.08\text{ g/mol}$ & 3 & 6 & 1 & $3 + 1.5 - 0.5 = \mathbf{4.0}$ & $0.50$ \\
Ácido Acético ($C_2H_4O_2$) & $60.05\text{ g/mol}$ & 2 & 4 & 2 & $2 + 1 - 1.0 = \mathbf{2.0}$ & $0.50$ \\
\bottomrule
\end{tabular}
\end{center}
\textit{Interpretación:} Observa cómo la acetona requiere más oxígeno ($\nu_{O_2} = 4.0$) que el ácido acético ($\nu_{O_2} = 2.0$), a pesar de tener una masa molar casi idéntica. Como consecuencia, la entalpía de combustión de la acetona ($\approx -1790\text{ kJ/mol}$) es más del doble de exotérmica que la del ácido acético ($\approx -874\text{ kJ/mol}$). ¡La variable $\nu_{O_2}$ contiene enorme poder predictivo!
\end{exercisebox}

\newpage

% ------------------------------------------------------------------------------
% CAPÍTULO 6: ANÁLISIS DEL TRABAJO PÚBLICO DEL DR. ARZOLA FLORES
% ------------------------------------------------------------------------------
\chapter{Análisis Especializado del Trabajo Público del Dr. Arzola Flores}

Karol Paola, en este capítulo dedicamos una sección específica y detallada a analizar con rigor científico la investigación pública desarrollada por el Dr. Jesús Andrés Arzola Flores y su grupo de colaboradores en la BUAP sobre la predicción de entalpías de combustión. Nuestro único y exclusivo objetivo pedagógico es aprender de su arquitectura y resultados para diseñar nuestro propio modelo predictivo limpio, transparente y autónomo.

\section{Ficha del Trabajo Público de Referencia}
\begin{itemize}[leftmargin=*]
    \item \textbf{Cita Científica Formal:} Saviñón-Flores, M. F., Arzola-Flores, J. A., García-Castro, M. A., Díaz-Sánchez, F., Vidal-Robles, E., \& Maruri-Valderrabano, F. A. (2025). \textit{Prediction of Standard Combustion Enthalpy of Organic Compounds Combining Machine Learning and Chemical Graph Theory: A Strategy}. \textbf{ACS Omega}, 10(36), 41828--41848.
    \item \textbf{Identificador DOI Oficial:} \url{https://doi.org/10.1021/acsomega.5c05927}
\end{itemize}

\section{¿Cuál fue el Problema Científico que Abordaron?}
El objetivo del grupo consistió en predecir de manera generalizada y precisa la entalpía estándar de combustión ($\Delta H_c^\circ$) de una gran biblioteca de moléculas orgánicas con alta diversidad estructural (hidrocarburos lineales, ramificados, ciclos alifáticos, aromáticos, heterociclos oxigenados y nitrogenados). Buscaban una alternativa cuantitativa a la experimentación destructiva en calorímetros de bomba, útil para evaluar rápidamente el potencial energético de nuevos compuestos y biocombustibles.

\section{El Dataset Utilizado y su Curaduría}
\begin{itemize}[leftmargin=*]
    \item \textbf{Tamaño Muestral:} \textbf{3,477 compuestos orgánicos curados}.
    \item \textbf{Fuentes de Datos:} Bases de datos termodinámicas de referencia internacional (NIST Chemistry WebBook, compilaciones termoquímicas calorimétricas históricas y reportes certificados).
    \item \textbf{Variable Objetivo ($y$):} Entalpía estándar de combustión a $298.15\text{ K}$ en $\text{kJ/mol}$ (con valores que abarcan desde moléculas pequeñas como el metano hasta hidrocarburos complejos de más de 30 carbonos, con entalpías entre $-800$ y $-16,000\text{ kJ/mol}$).
\end{itemize}

\section{¿Qué Crearon y Cómo lo Hicieron?}
El pipeline desarrollado por el grupo de investigación constó de las siguientes etapas secuenciales:

\begin{enumerate}[leftmargin=*]
    \item \textbf{Normalización de Identificadores:} Extracción y estandarización de cadenas SMILES canónicas para evitar representaciones redundantes de la misma estructura química.
    \item \textbf{Generación de la Matriz de Diseño ($\mathbf{X}$):} Extracción computacional de variables numéricas para cada una de las 3,477 moléculas.
    \item \textbf{Partición de Datos:} División del conjunto total en un 80\% de datos para entrenamiento (\textit{Training Set}) y un 20\% reservado estrictamente para prueba (\textit{Test Set}), asegurando reproducibilidad mediante semillas fijas.
    \item \textbf{Selección y Entrenamiento del Algoritmo Ganador:} Tras evaluar diversos estimadores (regresiones lineales y algoritmos de ensamble), el modelo que obtuvo el rendimiento más consistente y robusto fue un ensamble de árboles de decisión: \textbf{Random Forest Regressor}.
    \item \textbf{Afinamiento de Hiperparámetros:} Control riguroso del número de árboles (\texttt{n\_estimators}) y de la profundidad máxima (\texttt{max\_depth}) para mitigar el riesgo de sobreajuste (\textit{overfitting}).
\end{enumerate}

\section{Resultados Cuantitativos Alcanzados en su Publicación}
En el conjunto de prueba independiente (\textit{Test Set}), el modelo del Dr. Arzola Flores y colaboradores reportó los siguientes valores de referencia:
\begin{tcolorbox}[colback=boxbg,colframe=buapblue,title=Métricas Oficiales Reportadas en ACS Omega 2025]
\begin{align*}
    R^2 &\approx \mathbf{0.9810} \quad (\text{Coeficiente de determinación}) \\
    \text{MAE} &\approx \mathbf{287.5988\text{ kJ/mol}} \quad (\text{Error Absoluto Medio})
\end{align*}
\end{tcolorbox}

\section{Lecciones Didácticas para Nuestro Propio Modelo}
Analizar este trabajo nos deja lecciones invaluables para nuestro proyecto:
\begin{enumerate}[leftmargin=*]
    \item \textbf{Random Forest es el Estimador Ideal para este Target:} La relación entre la estructura molecular y la entalpía de combustión presenta no-linealidades inherentes (tensión de anillo en ciclos pequeños, efectos de resonancia aromática, interacciones estéricas). Los ensambles de árboles capturan estas particiones no lineales sin sufrir por problemas de escala.
    \item \textbf{La Métrica de Éxito es el MAE, no solo el $R^2$:} Como la base de datos cubre un rango numérico gigantesco (de $-890$ a $-16000\text{ kJ/mol}$), la varianza total es inmensa. Un $R^2 = 0.98$ puede sonar casi perfecto, pero un MAE de $287\text{ kJ/mol}$ representa el margen real de incertidumbre media del modelo.
    \item \textbf{Nuestro Camino Autónomo:} En lugar de complicarnos en esta etapa inicial con índices de grafos topológicos que aún no dominas, utilizaremos descriptores constitutivos, estequiométricos y fisicoquímicos directos (masa molar, balance de oxígeno, conteo de heteroátomos y enlaces). Diseñaremos nuestro propio pipeline claro, modular y 100\% comprensible para ti.
\end{enumerate}

\newpage

% ------------------------------------------------------------------------------
% CAPÍTULO 7: MACHINE LEARNING TEÓRICO Y MODELOS DE REGRESIÓN
% ------------------------------------------------------------------------------
\chapter{Machine Learning Fundamental: Algoritmos de Regresión Explicados}

Karol Paola, en este capítulo aprenderás qué hace internamente cada algoritmo de regresión, deduciendo sus ecuaciones y entendiendo por qué los árboles y ensambles son tan efectivos.

\section{¿Qué es un Problema de Regresión Supervisada?}
Disponemos de un conjunto de entrenamiento de $n$ moléculas:
\begin{equation}
    \mathcal{D} = \{(\mathbf{x}_1, y_1), (\mathbf{x}_2, y_2), \dots, (\mathbf{x}_n, y_n)\}
\end{equation}
donde $\mathbf{x}_i \in \mathbb{R}^p$ es el vector de variables moleculares y $y_i \in \mathbb{R}$ es la entalpía de combustión medida experimentalmente. Nuestro objetivo es aprender una función hipótesis $\hat{y} = f(\mathbf{x})$ tal que el error entre las predicciones $\hat{y}$ y los valores reales $y$ sea mínimo en moléculas nuevas que el modelo jamás haya visto.

\section{Regresión Lineal por Mínimos Cuadrados Ordinarios (OLS)}

\subsection{Ecuación del Modelo y Función de Costo (Loss)}
El modelo lineal asume que la entalpía predicha es una combinación lineal de las variables:
\begin{equation}
    \hat{y}_i = w_0 + w_1 x_{i1} + w_2 x_{i2} + \dots + w_p x_{ip} = \mathbf{w}^T \mathbf{x}_i + b
\end{equation}
Para encontrar los pesos óptimos $\mathbf{w}$, definimos la función de pérdida del \textbf{Error Cuadrático Medio (MSE)}:
\begin{equation}
    L(\mathbf{w}, b) = \frac{1}{n} \sum_{i=1}^n \left( y_i - (\mathbf{w}^T \mathbf{x}_i + b) \right)^2
\end{equation}

\subsection{Solución Analítica Matricial}
Expresando el sistema en notación matricial (añadiendo una columna de unos para el sesgo $b$ en la matriz $\mathbf{X} \in \mathbb{R}^{n \times (p+1)}$):
\begin{equation}
    L(\mathbf{w}) = \frac{1}{n} \|\mathbf{y} - \mathbf{X}\mathbf{w}\|^2 = \frac{1}{n} (\mathbf{y} - \mathbf{X}\mathbf{w})^T (\mathbf{y} - \mathbf{X}\mathbf{w})
\end{equation}
Derivando con respecto al vector de pesos $\mathbf{w}$ e igualando a cero:
\begin{equation}
    \nabla_{\mathbf{w}} L = -\frac{2}{n} \mathbf{X}^T (\mathbf{y} - \mathbf{X}\mathbf{w}) = \mathbf{0} \quad \Longrightarrow \quad \mathbf{X}^T \mathbf{X} \mathbf{w} = \mathbf{X}^T \mathbf{y}
\end{equation}
Si la matriz $\mathbf{X}^T \mathbf{X}$ es invertible, la solución analítica exacta es:
\begin{equation}
    \mathbf{w}^* = (\mathbf{X}^T \mathbf{X})^{-1} \mathbf{X}^T \mathbf{y}
\end{equation}

\section{Regularización: Evitando el Sobreajuste con Ridge y Lasso}

\subsection{El Dilema Sesgo-Varianza (\textit{Bias-Variance Tradeoff})}
\begin{itemize}[leftmargin=*]
    \item \textbf{Subajuste (\textit{Underfitting} / Alto Sesgo):} El modelo es demasiado simple (ej. asumir una línea recta cuando la relación es curva). Comete grandes errores tanto en entrenamiento como en prueba.
    \item \textbf{Sobreajuste (\textit{Overfitting} / Alta Varianza):} El modelo memoriza el ruido del entrenamiento. Tiene un error casi nulo en Train, pero falla estrepitosamente en Test.
\end{itemize}

\subsection{Regresión Ridge (Penalización $L_2$)}
Ridge añade a la función de pérdida una penalización cuadrática proporcional a la magnitud de los pesos:
\begin{equation}
    L_{\text{Ridge}}(\mathbf{w}) = \frac{1}{n} \sum_{i=1}^n (y_i - \mathbf{w}^T \mathbf{x}_i)^2 + \alpha \sum_{j=1}^p w_j^2 = \text{MSE} + \alpha \|\mathbf{w}\|_2^2
\end{equation}
donde $\alpha > 0$ es el hiperparámetro de regularización.  
\textbf{Efecto pedagógico:} Impide que ningún peso se vuelva astronómicamente grande, encogiendo los coeficientes hacia cero y aumentando la estabilidad del modelo.

\subsection{Regresión Lasso (Penalización $L_1$)}
Lasso penaliza el valor absoluto de los pesos:
\begin{equation}
    L_{\text{Lasso}}(\mathbf{w}) = \frac{1}{n} \sum_{i=1}^n (y_i - \mathbf{w}^T \mathbf{x}_i)^2 + \alpha \sum_{j=1}^p |w_j| = \text{MSE} + \alpha \|\mathbf{w}\|_1
\end{equation}
\textbf{Efecto pedagógico:} Debido a la geometría de las esquinas de la norma $L_1$, Lasso anula por completo los pesos de variables irrelevantes ($w_j = 0$), actuando como un seleccionador automático de variables químicas.

\section{Árboles de Decisión y Random Forest Regressor}

\subsection{¿Cómo Funciona un Árbol de Regresión Individual?}
Un árbol de decisión particiona el espacio de descriptores en regiones rectangulares hiperdimensionales $R_1, R_2, \dots, R_M$. Para cualquier molécula que caiga en la región $R_m$, la predicción es simplemente el promedio de las entalpías de las moléculas de entrenamiento de esa región:
\begin{equation}
    \hat{y}_{R_m} = \frac{1}{N_m} \sum_{i \in R_m} y_i
\end{equation}
En cada nodo, el algoritmo elige la variable $j$ y el punto de corte $s$ que maximizan la reducción de varianza (disminución del MSE).

\subsection{Random Forest: El Poder del Ensamble}
Un solo árbol de decisión profundo tiende a sobreajustar con facilidad. \textbf{Random Forest} resuelve esto promediando las predicciones de $B$ árboles descorrelacionados (típicamente $B = 100$ a $300$):
\begin{equation}
    \hat{y}_{\text{RF}}(\mathbf{x}) = \frac{1}{B} \sum_{b=1}^B T_b(\mathbf{x})
\end{equation}
Para asegurar que los árboles sean distintos entre sí, utiliza dos mecanismos aleatorios:
\begin{enumerate}[leftmargin=*]
    \item \textbf{Bootstrap Aggregating (Bagging):} Cada árbol se entrena con una muestra aleatoria del mismo tamaño $n$, tomada con reemplazo del dataset original.
    \item \textbf{Submuestreo de Características:} En cada división de cada nodo, el árbol solo puede elegir entre un subconjunto aleatorio de variables (típicamente $\sqrt{p}$ o $p/3$).
\end{enumerate}
\textit{Demostración estadística:} Si promediamos $B$ estimadores independientes, cada uno con varianza $\sigma^2$, la varianza del promedio es $\frac{\sigma^2}{B}$. ¡Random Forest reduce drásticamente la varianza del error sin aumentar el sesgo!

\section{Pipeline Maestro de Scikit-Learn (Sin Data Leakage)}

\begin{lstlisting}[language=Python]
# ==============================================================================
# PIPELINE MAESTRO DE MACHINE LEARNING EN INGENIERIA QUIMICA
# ==============================================================================
import numpy as np
import pandas as pd
from sklearn.model_selection import train_test_split
from sklearn.preprocessing import StandardScaler
from sklearn.pipeline import Pipeline
from sklearn.linear_model import LinearRegression, Ridge
from sklearn.ensemble import RandomForestRegressor
from sklearn.metrics import mean_absolute_error, mean_squared_error, r2_score

# 1. Separacion de matriz de features (X) y vector objetivo (y)
columnas_features = ['peso_mol', 'num_C', 'num_H', 'num_O', 'nu_O2', 'ratio_C_H']
X = df_limpio[columnas_features]
y = df_limpio['dH_combustion'] # Target en kJ/mol

# 2. Particion Train/Test estratificada en reproducibilidad (80% Train, 20% Test)
X_train, X_test, y_train, y_test = train_test_split(
    X, y, test_size=0.20, random_state=42
)

# 3. Configuracion de estimadores dentro de Pipelines para evitar fuga de datos
modelos = {
    "Regresion Lineal (OLS)": Pipeline([
        ('escalador', StandardScaler()),
        ('regresor', LinearRegression())
    ]),
    "Ridge (L2 regularizado)": Pipeline([
        ('escalador', StandardScaler()),
        ('regresor', Ridge(alpha=10.0))
    ]),
    "Random Forest (200 arboles)": RandomForestRegressor(
        n_estimators=200, max_depth=14, random_state=42, n_jobs=-1
    )
}

# 4. Entrenamiento y evaluacion comparativa
resultados = []
for nombre, modelo in modelos.items():
    modelo.fit(X_train, y_train) # Ajustar EXCLUSIVAMENTE con datos de Train
    
    pred_train = modelo.predict(X_train)
    pred_test = modelo.predict(X_test)
    
    resultados.append({
        "Algoritmo": nombre,
        "MAE Train (kJ/mol)": round(mean_absolute_error(y_train, pred_train), 2),
        "MAE Test (kJ/mol)": round(mean_absolute_error(y_test, pred_test), 2),
        "RMSE Test (kJ/mol)": round(np.sqrt(mean_squared_error(y_test, pred_test)), 2),
        "R2 Test": round(r2_score(y_test, pred_test), 4)
    })

df_eval = pd.DataFrame(resultados)
print(df_eval.to_string(index=False))
\end{lstlisting}

\newpage

% ------------------------------------------------------------------------------
% CAPÍTULO 8: VALIDACIÓN CIENTÍFICA Y DIAGNÓSTICO DE RESIDUALES
% ------------------------------------------------------------------------------
\chapter{Validación Científica, Métricas y Diagnóstico de Residuales}

Karol Paola, en ciencia no basta con decir que un modelo es bueno; debemos demostrar rigurosamente su precisión, sus límites de validez y sus posibles sesgos.

\section{Fórmulas Matemáticas de las Métricas de Error}

\begin{enumerate}[leftmargin=*]
    \item \textbf{Error Absoluto Medio (MAE - La Métrica Reina):}
    \begin{equation}
        \text{MAE} = \frac{1}{n_{\text{test}}} \sum_{i=1}^{n_{\text{test}}} |y_i - \hat{y}_i|
    \end{equation}
    Mide el error promedio en las unidades físicas reales del fenómeno ($\text{kJ/mol}$). Un MAE de $280\text{ kJ/mol}$ indica que, en promedio, la predicción del modelo se desvía en $\pm 280\text{ kJ/mol}$ del valor calorimétrico experimental.
    \item \textbf{Raíz del Error Cuadrático Medio (RMSE):}
    \begin{equation}
        \text{RMSE} = \sqrt{\frac{1}{n_{\text{test}}} \sum_{i=1}^{n_{\text{test}}} (y_i - \hat{y}_i)^2}
    \end{equation}
    Al elevar las diferencias al cuadrado antes de promediar, el RMSE penaliza severamente los errores grandes o atípicos (\textit{outliers}). Si el RMSE es significativamente mayor que el MAE, significa que el modelo comete fallos extremos en un subgrupo de moléculas.
    \item \textbf{Coeficiente de Determinación ($R^2$):}
    \begin{equation}
        R^2 = 1 - \frac{\sum_{i=1}^{n_{\text{test}}} (y_i - \hat{y}_i)^2}{\sum_{i=1}^{n_{\text{test}}} (y_i - \bar{y})^2} = 1 - \frac{SS_{\text{res}}}{SS_{\text{tot}}}
    \end{equation}
    Representa la fracción de la varianza del target explicada por el modelo predictivo.
\end{enumerate}

\begin{warningbox}{¡Cuidado con el Espejismo del $R^2$ en Termoquímica!}
Como la entalpía de combustión de compuestos orgánicos abarca un rango enorme (desde $-890\text{ kJ/mol}$ para el metano hasta más de $-15000\text{ kJ/mol}$ para alcanos pesados), el denominador de la varianza total ($SS_{\text{tot}}$) es colosal. Un modelo deficiente que se equivoque por $\pm 400\text{ kJ/mol}$ en cada predicción puede arrojar tranquilamente un $R^2 = 0.975$. \textbf{Tu criterio de calidad en laboratorio debe ser siempre el MAE}.
\end{warningbox}

\section{El Diagnóstico Gráfico de Residuales ($e_i = y_i - \hat{y}_i$)}
Grafica siempre los residuales frente a las entalpías predichas:
\begin{lstlisting}[language=Python]
import matplotlib.pyplot as plt

y_pred = modelos["Random Forest (200 arboles)"].predict(X_test)
residuales = y_test - y_pred

plt.figure(figsize=(7.5, 4.8))
plt.scatter(y_pred, residuales, alpha=0.5, color='navy', edgecolors='k', s=35)
plt.axhline(0, color='red', linestyle='--', linewidth=1.5)
plt.title('Diagnostico de Residuales: Random Forest Regressor', fontsize=12)
plt.xlabel('$\Delta H_c^\circ$ Predicho (kJ/mol)')
plt.ylabel('Residual ($y_{\mathrm{exp}} - y_{\mathrm{pred}}$) [kJ/mol]')
plt.grid(True, linestyle=':', alpha=0.6)
plt.tight_layout()
plt.show()
\end{lstlisting}

\noindent\textbf{Cómo interpretar la gráfica de residuales como una profesional:}
\begin{itemize}[leftmargin=*]
    \item \textbf{Banda aleatoria horizontal (Homocedasticidad):} La nube de puntos se distribuye simétricamente alrededor de la línea cero sin ensancharse. \textit{El modelo es robusto e insesgado.}
    \item \textbf{Forma de embudo (\textit{Funnel shape} / Heterocedasticidad):} La nube es estrecha a entalpías pequeñas y se ensancha a entalpías muy negativas. \textit{Indica que la incertidumbre crece en moléculas grandes.}
    \item \textbf{Puntos solitarios muy alejados ($|e_i| > 1500\text{ kJ/mol}$):} Son anomalías que requieren inspección química (¿SMILES mal escrito?, ¿error tipográfico en el CSV original?).
\end{itemize}

\newpage

% ------------------------------------------------------------------------------
% CAPÍTULO 9: METAS DEL VIERNES Y HOJA DE RUTA HASTA DICIEMBRE 2026
% ------------------------------------------------------------------------------
\chapter{Plan de Trabajo: Metas del Viernes y Hoja de Ruta a Diciembre}

Karol Paola, este capítulo define con exactitud qué debes alcanzar al concluir tu primera semana de trabajo y cómo se estructura el camino cronológico hasta el hito del 10 de diciembre de 2026.

\section{Checklist Obligatoria para este Primer Viernes}
Antes de cerrar tu sesión de trabajo de este viernes, debes verificar el cumplimiento estricto de cada uno de los siguientes puntos:

\begin{protocolbox}{Checklist Operativa del Primer Viernes}
\begin{itemize}[leftmargin=*]
    \item[$\square$] \textbf{Entorno PyCharm Funcional:} El entorno virtual \texttt{.venv} está configurado y el intérprete ejecuta Python 3.10+ sin errores de importación.
    \item[$\square$] \textbf{Estructura Limpia de Carpetas:} El archivo \texttt{base\_entalpias.csv} está colocado en \texttt{data/raw/} y no disperso en el escritorio o en carpetas temporales.
    \item[$\square$] \textbf{Cuaderno Inicial Creado:} El archivo interactivo \texttt{notebooks/01\_exploracion\_inicial.ipynb} está conectado al kernel de \texttt{.venv} y carga la base de datos con rutas relativas (\texttt{os.path.join}).
    \item[$\square$] \textbf{Auditoría de Datos Realizada:} Se han ejecutado \texttt{df.info()} y \texttt{df.describe()}, identificando el número total de filas, la columna target (\texttt{dH\_combustion}) y confirmando que las unidades son $\text{kJ/mol}$.
    \item[$\square$] \textbf{Detección de Nulos y Duplicados:} Se han cuantificado los valores nulos con \texttt{df.isnull().sum()} y los duplicados con \texttt{df.duplicated(subset=['SMILES']).sum()}.
    \item[$\square$] \textbf{ChemPy Operativo:} Se ha probado en un script el balanceo estequiométrico automático de la combustión del metano y etanol con \texttt{chempy.chemistry.balance\_stoichiometry}.
    \item[$\square$] \textbf{Git y GitHub Sincronizados:} El archivo \texttt{.gitignore} excluye correctamente \texttt{.venv/}, se ha realizado el commit descriptivo de la semana y el comando \texttt{git push origin main} ejecutó limpiamente.
\end{itemize}
\end{protocolbox}

\section{Hoja de Ruta Maestra por Fases (Septiembre a Diciembre 2026)}

\begin{table}[h!]
\centering
\small
\begin{tabularx}{\textwidth}{c p{3.2cm} p{4.2cm} X}
\toprule
\textbf{Fase} & \textbf{Periodo} & \textbf{Actividad Central} & \textbf{Entregable Oficial de Cierre} \\
\midrule
\rowcolor{codebg}
\textbf{F0-ENV} & Sep 04 -- Sep 05 & Setup de entorno, ChemPy y Git. & \texttt{F0-ENV\_ACT01\_CFG\_entorno\_pycharm\_v1.0\_FINAL.txt} \newline \texttt{F0-ENV\_ACT02\_SCR\_funciones\_chempy\_v1.0\_FINAL.py} \\
\addlinespace
\textbf{F1-DAT} & Sep 05 -- Sep 12 & Curaduría de datos, unidades y nulos. & \texttt{F1-DAT\_ACT03\_DAT\_dataset\_curado\_v1.0\_FINAL.csv} \\
\rowcolor{codebg}
\textbf{F2-EDA} & Sep 12 -- Sep 26 & Detección de outliers (IQR) y cálculo de descriptores moleculares. & \texttt{F2-EDA\_ACT04\_NBK\_exploracion\_distribuciones\_v1.0\_FINAL.ipynb} \newline \texttt{F2-EDA\_ACT05\_SCR\_calculo\_descriptores\_moleculares\_v1.0\_FINAL.py} \\
\addlinespace
\textbf{F3-MLB} & Sep 26 -- Oct 10 & Modelos lineales (OLS, Ridge) y ajuste de Random Forest Regressor. & \texttt{F3-MLB\_ACT06\_NBK\_modelos\_lineales\_baseline\_v1.0\_FINAL.ipynb} \newline \texttt{F3-MLB\_ACT07\_SCR\_random\_forest\_regressor\_v1.0\_FINAL.py} \\
\rowcolor{codebg}
\textbf{F4-VAL} & Oct 10 -- Oct 17 & Validación cruzada 10-fold y análisis gráfico de residuales. & \texttt{F4-VAL\_ACT08\_NBK\_diagnostico\_residuales\_v1.0\_FINAL.ipynb} \\
\addlinespace
\textbf{F5-REP} & Oct 17 -- Oct 31 & Benchmarking y comparación de métricas contra literatura. & \texttt{F5-REP\_ACT09\_REP\_informe\_benchmarking\_literatura\_v1.0\_FINAL.pdf} \\
\rowcolor{codebg}
\textbf{F6-SIM} & Oct 31 -- Nov 14 & Simulación termodinámica de combustión en DWSIM. & \texttt{F6-SIM\_SIM01\_SIM\_reactor\_combustion\_dwsim\_v1.0\_FINAL.dwsim} \\
\addlinespace
\textbf{F7-EQK} & Nov 14 -- Dic 10 & Puente al equilibrio químico ($\Delta G^\circ, K_{eq}$) y entrega final. & \texttt{F7-EQK\_ACT10\_SCR\_calculador\_gibbs\_keq\_v1.0\_FINAL.py} \newline \texttt{F7-EQK\_DOC01\_REP\_manuscrito\_final\_propuesta\_v1.0\_FINAL.pdf} \\
\bottomrule
\end{tabularx}
\caption{Cronograma integral del proyecto hacia el 10 de diciembre de 2026.}
\end{table}

\section{El Hito Decisivo del 10 de Diciembre de 2026}
¿Por qué todo este esfuerzo de otoño tiene como horizonte el 10 de diciembre de 2026? Porque el modelado de \textbf{Equilibrio Químico} en sistemas complejos multicomponente (reacciones simultáneas, balance de energía libre de Gibbs, coeficientes de actividad no ideales y cálculo de constantes $K_{eq}$) requiere que domines con total naturalidad el manejo de datos moleculares, el cálculo estequiométrico automatizado y los pipelines de regresión. Las entalpías de combustión son tu plataforma de despegue; el equilibrio químico será tu consagración como investigadora autónoma.

\newpage

% ------------------------------------------------------------------------------
% CAPÍTULO 10: CATÁLOGO DE ERRORES FRECUENTES Y SOLUCIONES
% ------------------------------------------------------------------------------
\chapter{Catálogo de Errores Frecuentes y Cómo Resolverlos}

Aprender a programar en ciencia de datos es, en gran medida, aprender a diagnosticar y corregir errores sin frustrarse. A continuación se presentan los problemas más comunes y su solución directa:

\begin{enumerate}[leftmargin=*]
    \item \textbf{\texttt{FileNotFoundError} al cargar el CSV:}  
    \textit{Causa:} El script se ejecuta desde una carpeta distinta a la que esperabas o utilizaste una ruta absoluta.  
    \textit{Solución:} Verifica tu directorio de trabajo actual con \texttt{os.getcwd()} y construye la ruta relativa con \texttt{os.path.join("..", "data", "raw", "base\_entalpias.csv")}.
    
    \item \textbf{\texttt{ModuleNotFoundError: No module named 'chempy'}:}  
    \textit{Causa:} Instalaste la librería en el Python global del sistema y no en el entorno virtual \texttt{.venv} del proyecto.  
    \textit{Solución:} Abre la terminal en PyCharm, verifica que comience con \texttt{(.venv)} y corre \texttt{pip install chempy}.
    
    \item \textbf{Fuga de Datos (\textit{Data Leakage}):}  
    \textit{Causa:} Aplicar \texttt{StandardScaler().fit()} sobre todo el dataset antes de hacer \texttt{train\_test\_split}.  
    \textit{Solución:} Emplea siempre un objeto \texttt{Pipeline([('scaler', StandardScaler()), ('model', Ridge())])}. El pipeline ajusta el escalador únicamente sobre el Train Set.
    
    \item \textbf{Inconsistencia de Unidades Termodinámicas:}  
    \textit{Causa:} Combinar columnas reportadas en $\text{kcal/mol}$ con datos en $\text{kJ/mol}$.  
    \textit{Solución:} Aplica el factor de conversión estricto ($1\text{ kcal} = 4.184\text{ kJ}$) al momento de cargar el dataset crudo.
    
    \item \textbf{Sobreajuste Severo en Random Forest (\textit{Overfitting}):}  
    \textit{Causa:} Dejar \texttt{max\_depth=None}, permitiendo que cada árbol crezca hasta que cada hoja tenga una sola molécula (memorización pura).  
    \textit{Solución:} Acota la profundidad fijando \texttt{max\_depth=12} o \texttt{max\_depth=14} y aumenta \texttt{min\_samples\_split=5}.
    
    \item \textbf{Subir el Entorno Virtual a GitHub:}  
    \textit{Causa:} Olvidar agregar \texttt{.venv/} al archivo \texttt{.gitignore}.  
    \textit{Solución:} Añade \texttt{.venv/} en una línea de \texttt{.gitignore} y ejecuta \texttt{git rm -r --cached .venv} para desrastrearlo sin borrarlo de tu disco.
    
    \item \textbf{Resultados no reproducibles entre corridas:}  
    \textit{Causa:} No fijar la semilla aleatoria en los algoritmos estocásticos.  
    \textit{Solución:} Especifica siempre \texttt{random\_state=42} tanto en \texttt{train\_test\_split} como en \texttt{RandomForestRegressor}.
\end{enumerate}

\vfill
\begin{center}
    \rule{0.6\textwidth}{0.6pt}\\[0.2cm]
    {\small \textbf{Manual Didáctico de Formación Progresiva en Termoquímica y Machine Learning}}\\[0.1cm]
    {\footnotesize Benemérita Universidad Autónoma de Puebla $\cdot$ Facultad de Ingeniería Química}
\end{center}

\end{document}
